\begin{examplebox}
	\textbf{\textcolor{CExample}{例 1（基础）}}
	\textbf{\textcolor{CExample}{题目：}}
	计算极限 $\displaystyle\lim_{x\to 0}\frac{\sin x}{x}$．
	\textbf{\textcolor{CExample}{解题步骤：}}
	\begin{description}[style=nextline, leftmargin=2.8em, labelsep=0.5em, itemsep=0.45em]
		\item[\textbf{第一步（判断类型）：}]
			当 $x\to0$ 时，分子 $\sin x\to0$，分母 $x\to0$，故极限为 $\frac{0}{0}$ 型未定式，满足洛必达法则条件．
			\par\textbf{理由：} 分子分母可导，且分母导数 $1\neq0$，极限为 $\frac{0}{0}$．
		\item[\textbf{第二步（求导计算）：}]
			对分子分母分别求导，得 $\lim_{x\to0}\frac{\cos x}{1}$．
			\par\textbf{理由：} 洛必达法则：原极限等于导数比极限，前提是导数比极限存在．
		\item[\textbf{第三步（求取极限）：}]
			$\lim_{x\to0}\frac{\cos x}{1}=1$，故原极限为 $1$．
			\par\textbf{理由：} $\cos0=1$．
	\end{description}
	\textbf{\textcolor{CExample}{关键结论：}}
	\textbf{$\displaystyle\lim_{x\to 0}\frac{\sin x}{x}=1$．}

	\medskip
	\textbf{\textcolor{CExample}{例 2（稍变形）}}
	\textbf{\textcolor{CExample}{题目：}}
	已知 $a\le \frac{e^x-1}{x}$ 对 $x>0$ 恒成立，求实数 $a$ 的取值范围．
	\textbf{\textcolor{CExample}{解题步骤：}}
	\begin{description}[style=nextline, leftmargin=2.8em, labelsep=0.5em, itemsep=0.45em]
		\item[\textbf{第一步（分离参数）：}]
			不等式化为 $a\le h(x)$，其中 $h(x)=\frac{e^x-1}{x}$．
			\par\textbf{理由：} 分离参数是处理恒成立问题的基本方法．
		\item[\textbf{第二步（识别未定式）：}]
			当 $x\to0^+$ 时，$e^x-1\to0$，$x\to0$，故 $h(x)$ 在端点 $0$ 处极限为 $\frac{0}{0}$ 型未定式．
			\par\textbf{理由：} 端点极限为未定式时，可利用洛必达法则求出该极限．
		\item[\textbf{第三步（洛必达计算）：}]
			$\displaystyle\lim_{x\to0^+}h(x)=\lim_{x\to0^+}\frac{e^x}{1}=1$．
			\par\textbf{理由：} 分子分母分别求导，分母导数 $1\neq0$，导数比极限存在．
		\item[\textbf{第四步（必要条件）：}]
			由 $a\le h(x)$ 恒成立，得 $a\le \lim_{x\to0^+}h(x)=1$，因此 $a\le1$．
			\par\textbf{理由：} 若 $a$ 大于该极限，则在 $x$ 足够接近 $0$ 时不等式不成立．这是必要条件，充分性需另行验证．
	\end{description}
	\textbf{\textcolor{CExample}{关键结论：}}
	\textbf{$a$ 的取值范围是 $(-\infty,1]$（必要，需验证充分性）．}

	\medskip
	\textbf{\textcolor{CExample}{例 3（综合应用）}}
	\textbf{\textcolor{CExample}{题目：}}
	已知 $a\ge \frac{\ln x}{x-1}$ 对 $x>1$ 恒成立，求实数 $a$ 的取值范围．
	\textbf{\textcolor{CExample}{解题步骤：}}
	\begin{description}[style=nextline, leftmargin=2.8em, labelsep=0.5em, itemsep=0.45em]
		\item[\textbf{第一步（分离参数）：}]
			不等式化为 $a\ge h(x)$，其中 $h(x)=\frac{\ln x}{x-1}$．
			\par\textbf{理由：} 分离参数将恒成立问题转化为函数最值问题．
		\item[\textbf{第二步（识别未定式）：}]
			当 $x\to1^+$ 时，$\ln x\to0$，$x-1\to0$，故 $h(x)$ 在端点 $1$ 处极限为 $\frac{0}{0}$ 型未定式．
			\par\textbf{理由：} 端点出现未定式，可应用洛必达法则．
		\item[\textbf{第三步（洛必达计算）：}]
			$\displaystyle\lim_{x\to1^+}h(x)=\lim_{x\to1^+}\frac{1/x}{1}=1$．
			\par\textbf{理由：} 分子分母分别求导，导数比极限存在．
		\item[\textbf{第四步（必要范围）：}]
			由 $a\ge h(x)$ 恒成立，得 $a\ge \lim_{x\to1^+}h(x)=1$，即 $a\ge1$．这是必要条件．
			\par\textbf{理由：} 若 $a$ 小于该极限，则 $x$ 充分接近 $1$ 时不等式不成立．
		\item[\textbf{第五步（充分性提示）：}]
			可证明 $h(x)\le 1$（例如构造函数 $\phi(x)=\ln x-(x-1)\le0$），所以当 $a\ge1$ 时不等式恒成立．因此 $a\ge1$ 是充要条件．
			\par\textbf{理由：} 充分性通过函数最值或不等式证明．
	\end{description}
	\textbf{\textcolor{CExample}{关键结论：}}
	\textbf{$a$ 的取值范围是 $[1,+\infty)$（充要）．}
\end{examplebox}
